\chapter{Euler Equation (Fluid)}

\section*{Introduction}

Imagine a swarm of bees flying through air. Each bee (fluid parcel) is pushed by pressure differences (high to low pressure) and pulled by gravity. The Euler equation says: the rate at which a bee accelerates equals the pressure gradient force divided by its mass density, plus gravity. In a real fluid, there would be friction between neighboring bees (viscosity), but the Euler equations ignore this---they describe the ``skeleton'' of fluid motion, to which viscosity adds the ``flesh.''


For an inviscid fluid with velocity $\mathbf{v}$, pressure $P$, density $\rho$, and gravity $\mathbf{g}$:
\[
\rho\left(\frac{\partial \mathbf{v}}{\partial t} + \mathbf{v} \cdot \nabla \mathbf{v}\right) = -\nabla P + \rho \mathbf{g}
\]
together with the continuity equation $\partial \rho / \partial t + \nabla \cdot (\rho \mathbf{v}) = 0$.

\section*{Fundamental Equations Used}

The derivation starts from Newton's second law applied to an infinitesimal fluid parcel.

\section*{Assumptions}

\textbf{Assumptions:} The fluid is inviscid (no internal friction). The flow is a continuum. Body forces are conservative.


\textbf{Limitations:} Cannot describe boundary layers, skin friction, or viscous-dominated flows. Fails for turbulence at high Reynolds number. Does not account for surface tension or thermal conduction (unless added separately).


\section*{Mathematical Derivation}

\textbf{Step 1: Start from Newton's second law for a fluid parcel.}

Consider a small fluid parcel of volume $\delta V$ and mass $\delta m = \rho\,\delta V$. Newton's second law applied to this parcel gives:
\begin{equation}
\delta m \cdot \frac{D\mathbf{v}}{Dt} = \delta\mathbf{F}_{\text{pressure}} + \delta\mathbf{F}_{\text{body}}
\end{equation}
where $D\mathbf{v}/Dt$ is the material (Lagrangian) derivative of the velocity, following the parcel.

\textbf{Step 2: Express the pressure force.}

The pressure force on the parcel is the integral of $-P\,\hat{\mathbf{n}}$ over its surface $S$. By the divergence theorem:
\begin{equation}
\delta\mathbf{F}_{\text{pressure}} = -\iint_S P\,d\mathbf{S} = -\iiint_{\delta V}\nabla P\,dV = -\nabla P\,\delta V
\end{equation}

\textbf{Step 3: Express the body force.}

For a conservative body force per unit mass $\mathbf{g}$ (e.g., gravity), the total body force is:
\begin{equation}
\delta\mathbf{F}_{\text{body}} = \rho\,\mathbf{g}\,\delta V
\end{equation}

\textbf{Step 4: Divide by the volume.}

Substituting into Newton's law and dividing by $\delta V$:
\begin{equation}
\rho\frac{D\mathbf{v}}{Dt} = -\nabla P + \rho\mathbf{g}
\end{equation}

\textbf{Step 5: Expand the material derivative.}

The material derivative is:
\begin{equation}
\frac{D\mathbf{v}}{Dt} = \frac{\partial \mathbf{v}}{\partial t} + (\mathbf{v} \cdot \nabla)\mathbf{v}
\end{equation}
The first term is the local acceleration (change at a fixed point), and the second is the convective acceleration (change due to the parcel moving into a region of different velocity).

\textbf{Step 6: Write the Euler equation.}

Combining everything:
\begin{equation}
\boxed{\rho\left(\frac{\partial \mathbf{v}}{\partial t} + \mathbf{v}\cdot\nabla\mathbf{v}\right) = -\nabla P + \rho\mathbf{g}}
\end{equation}
together with the continuity equation $\partial\rho/\partial t + \nabla\cdot(\rho\mathbf{v}) = 0$ and an equation of state (e.g., $P = P(\rho, T)$) that closes the system. The equation is first-order in time and nonlinear due to the convective term $\mathbf{v}\cdot\nabla\mathbf{v}$. The Navier--Stokes equation adds the viscous term $\mu\nabla^2\mathbf{v}$ to the right side; the Euler equation is recovered in the limit of zero viscosity.

\section*{Characteristics of the Equation}

\begin{itemize}
\item \textbf{Nonlinear:} The convective term $\mathbf{v} \cdot \nabla \mathbf{v}$ makes the equation nonlinear---solutions can develop shocks from smooth data.
\item \textbf{Reversible:} Without viscosity, the equations are time-reversal symmetric---reversing velocities makes the fluid retrace its path.
\item \textbf{Bernoulli's equation:} For steady, irrotational flow: $P + \frac{1}{2}\rho v^2 + \rho gh = \text{const}$.
\item \textbf{Vorticity:} Vortex lines are material lines (Kelvin's circulation theorem): circulation is conserved.
\item \textbf{Sonic transitions:} For compressible flow, shock waves form at subsonic--supersonic transitions.
\end{itemize}
\section*{Solution Methods}

\begin{enumerate}
\item \textbf{Bernoulli integration:} For steady, incompressible, irrotational flow, integrate along streamlines.
\item \textbf{Potential flow:} For irrotational flow, $\mathbf{v} = \nabla\phi$, and equations reduce to Laplace's equation $\nabla^2 \phi = 0$.
\item \textbf{Method of characteristics:} For 1D compressible flow, reduce to ODEs along characteristic curves---useful for shock analysis.
\item \textbf{Numerical methods:} Finite volume, finite element, and spectral methods for complex geometries.
\end{enumerate}
\section*{Verifications}

\begin{itemize}
\item \textbf{Bernoulli experiments:} Measured pressure--velocity relationships in Venturi tubes match predictions.
\item \textbf{Potential flow:} Flow around a cylinder matches measurements outside the boundary layer.
\item \textbf{Sound speed:} Predicted $c = \sqrt{\gamma P / \rho}$ matches experimental values.
\item \textbf{Shock conditions:} Euler equations correctly predict jump conditions across shocks (Rankine--Hugoniot).
\end{itemize}
\section*{Applications}

\begin{itemize}
\item \textbf{Aerodynamics:} Inviscid theory (panel methods, conformal mapping) predicts lift (Kutta--Joukowski theorem).
\item \textbf{Weather:} Large-scale atmospheric dynamics is approximately inviscid---Euler equations with Coriolis force govern geostrophic flow.
\item \textbf{Ocean currents:} Large-scale circulation governed by Euler equations modified by rotation and stratification.
\item \textbf{Acoustics:} Sound waves are small-amplitude perturbations linearized about a uniform state.
\item \textbf{Rocket propulsion:} Compressible flow through nozzles predicts thrust and exit conditions.
\end{itemize}
\section*{What Richard Feynman Would Have Asked}

\subsection*{1. Plain-English Translation}
\textit{``What is the Euler equation for fluids actually saying?''}

It says: ``A fluid accelerates because of pressure differences and gravity---nothing else.'' Push on a fluid from the high-pressure side, it moves toward low pressure. Gravity pulls it down. That is the entire content. The nonlinear term $\mathbf{v} \cdot \nabla \mathbf{v}$ is just the statement that velocity changes as the fluid moves from one place to another---the ``convective acceleration.''

The deepest implication: in a frictionless fluid, energy is conserved---no dissipation. The fluid can slosh back and forth forever without losing energy. This idealization captures the essential physics of most flows: pressure pushes, gravity pulls, and inertia carries.

\subsection*{2. Localized Mechanics}
\textit{``What is happening to a single fluid parcel at a given instant?''}

At any instant, a fluid parcel has a certain velocity and pressure. The equation says the parcel accelerates in the direction of the pressure gradient (high to low) and in the direction of gravity. The acceleration is the material derivative $D\mathbf{v}/Dt = \partial\mathbf{v}/\partial t + \mathbf{v} \cdot \nabla\mathbf{v}$, including both local change and convective transport.

He would press you on the convective term: $\mathbf{v} \cdot \nabla\mathbf{v}$ is not a force---it is kinematic. It says a parcel moving into a region of different velocity changes to match the local flow. This is like a car accelerating when entering a faster lane---not because a force pushed it, but because it moved where the ``speed limit'' is different.

\subsection*{3. Testing Extreme Limits}
\textit{``What happens at extreme speeds or densities?''}

At very low speeds (creeping flow, $Re \ll 1$), viscous forces dominate and Euler is useless. At very high speeds (hypersonic, $M > 5$), compressibility dominates---Euler still applies, but density changes dramatically across shocks. At the speed of sound ($M = 1$), flow transitions from subsonic to supersonic, and Euler predicts sonic lines and shock waves.

In extreme density (neutron star interiors, early universe), Euler must be replaced by relativistic fluid dynamics, where the stress--energy tensor replaces Newtonian pressure and density.

\subsection*{4. Dimensionality and Geometry}
\textit{``How does the geometry of the flow domain affect the Euler equations?''}

In 1D, they reduce to hyperbolic PDEs supporting shock waves (Riemann problem). In 2D, they support vortices, stagnation points, and conformal mapping solutions. In 3D, helical vortices, vortex rings, and complex structures emerge.

Boundary geometry shapes the flow: a cylinder gives symmetric streamlines with no drag (d'Alembert's paradox); an airfoil with the Kutta condition produces lift. The geometry is the boundary condition that shapes the solution.

\subsection*{5. Cross-Disciplinary Connection}
\textit{``Where else does this same mathematical structure appear?''}

The Euler equations have the same form as equations for any continuum: in plasma physics (Lorentz force replaces gravity), magnetohydrodynamics (magnetic forces added), traffic flow (vehicle density and velocity), and crowd dynamics (pedestrian density). The mathematical structure---inertial acceleration equals force density---is universal for any continuous medium.


%=============================================================
% Chapter 24: Euler-Bernoulli Beam Equation
%=============================================================
\section*{FAQ}

\textbf{Q1: If the Euler equations ignore viscosity, why are they useful?}

A: Because for most of the flow field (away from walls), viscous effects are negligible at high Reynolds number. The Euler equations capture the dominant dynamics---pressure, inertia, gravity---while viscosity acts only in thin boundary layers. This ``outer flow'' approximation starts all aerodynamic analysis.

\textbf{Q2: What is the difference between Euler and Navier--Stokes?}

A: Navier--Stokes adds viscosity: the extra term $\mu \nabla^2 \mathbf{v}$ accounts for internal friction. At high Reynolds number, viscosity matters only near surfaces (boundary layers), so Euler is good for bulk flow. At low Reynolds number, viscosity dominates everywhere, and full Navier--Stokes is needed.

\textbf{Q3: Can the Euler equations predict turbulence?}

A: No. Turbulence requires viscosity to generate small-scale structure. The Euler equations are deterministic and time-reversible---they cannot produce chaotic, dissipative turbulence. However, they can predict instabilities that lead to turbulence (linear stability analysis).
\section*{Important Bibliography}

\begin{enumerate}
\item Batchelor, G.K., \textit{An Introduction to Fluid Dynamics} (Cambridge University Press, 1967), Chapter 3.
\item Kundu, P.K., Cohen, I.M., and Dowling, D.R., \textit{Fluid Mechanics}, 7th ed. (Academic Press, 2016), Chapter 5.
\item Landau, L.D. and Lifshitz, E.M., \textit{Fluid Mechanics}, 2nd ed. (Butterworth-Heinemann, 1987), Chapter 2.
\item Anderson, J.D., \textit{Modern Compressible Flow: With Historical Perspective}, 3rd ed. (McGraw-Hill, 2003), Chapter 2.
\end{enumerate}

